\section{Peterson's deformations of quadrics}\label{sec:pet}

Although Peterson \cite{P} discusses all types of quadrics in the
complexif\/ied Euclidean space
\[
\big(\mathbb{C}^3,\langle \cdot,\cdot\rangle\big),\qquad \langle x,y\rangle:=x^Ty,\qquad |x|^2:=x^Tx\quad \mathrm{for}\ \
x,y\in\mathbb{C}^3
\]
and their totally real cases, we shall only discuss quadrics of
the type $\sum\limits_{j=0}^2\frac{(x_j)^2}{a_j}=1$, $a_j\in\mathbb{C}^*$
distinct, since the remaining cases of quadrics should follow by
similar computations. Their totally real cases (that is
$(x_j)^2$, $a_j\in\mathbb{R}$) are discussed in detail in Peterson~\cite{P}, so we shall not insist on this aspect.

\begin{remark}
It is less known since the classical times that there are many
types of quadrics from a~complex metric point of view, each coming
with its own totally real cases (real valued (in)def\/inite linear
element); among these quadrics for example the quadric
$(x_0-ix_1)x_2-(x_0+ix_1)=0$ is rigidly {\it applicable}
(isometric) to all quadrics of its confocal family and to all its
homothetic quadrics. It is Peterson who f\/irst introduced the idea
of {\it ideal applicability} (for example a real surface may be
applicable to a totally real space-like surface
$\subset\mathbb{R}^2\times(i\mathbb{R})$ of a~complexif\/ied real
ellipsoid, so it is ideally applicable on the real ellipsoid).
\end{remark}

With $\{e_j\}_{j=0,1,2}$, $e_j^Te_k=\delta_{jk}$ the standard basis
of $\mathbb{C}^3$  and the functions $f=f(z,u^1)$, $g=g(z,u^1)$,
$h=h(z,u^2)$ depending on the parameter(s) $z=(z_1,z_2,\dots )$ to be
determined later we have the surfaces
\begin{gather}
\mathcal{X}_z:=\cos\big(u^2\big)f\big(z,u^1\big)\big(\cos\big(g\big(z,u^1\big)\big)e_0+\sin
\big(g\big(z,u^1\big)\big)e_1\big)+h\big(z,u^2\big)e_2.
\label{eq:surf}
\end{gather}
Note that the f\/ields
$\partial_{u^1}\mathcal{X}_z|_{u^1=\mathrm{const}}$,
$\partial_{u^2}\mathcal{X}_z|_{u^2=\mathrm{const}}$ generate developables
(cylinders with gene\-ra\-tors perpendicular on the third axis and
cones centered on the third axis), so $(u^1,u^2)$ is a~conjugate
system on $\mathcal{X}_z$ for every $z$; in fact all surfaces have
conjugate systems arising this way and can be parameterized as
\[
x=f\big(u^1,u^2\big)\big(\cos\big(u^1\big)e_0+\sin\big(u^1\big)e_1\big)+g\big(u^1,u^2\big)e_2, \qquad
\partial_{u^1}\left(\partial_{u^2}\left(\frac{g}{f}\right){\Big/}\partial_{u^2}\left(\frac{1}{f}\right)\right)=0.
\]

The quadric $\sum\limits_{j=0}^2\frac{(x_j)^2}{a_j}=1$ is parameterized
by the spherical coordinates
\[
\mathcal{X}=\sqrt{a_0}\cos\big(u^2\big)\cos\big(u^1\big)e_0
+\sqrt{a_1}\cos\big(u^2\big)\sin\big(u^1\big)e_1+\sqrt{a_2}\sin\big(u^2\big)e_2.
\]
We have
\begin{gather*}
|d\mathcal{X}_z|^2=\cos^2\big(u^2\big)\big(f'^2\big(z,u^1\big)
+f^2\big(z,u^1\big)g'^2\big(z,u^1\big)\big)\big(du^1\big)^2
\\
\phantom{|d\mathcal{X}_z|^2=}{}
+\tfrac{1}{2}d\big(\cos^2\big(u^2\big)\big)d\big(f^2\big(z,u^1\big)\big)
+\big(f^2\big(z,u^1\big)\sin^2\big(u^2\big)+h'^2\big(z,u^2\big)\big)\big(du^2\big)^2,
\\
|d\mathcal{X}|^2
=\cos^2\big(u^2\big)\big(a_1-(a_1-a_0)\sin^2\big(u^1\big)\big)\big(du^1\big)^2\\
\phantom{|d\mathcal{X}|^2=}{}
+\tfrac{1}{2}d\big(\cos^2\big(u^2\big)\big)
d\big(a_0+(a_1-a_0)\sin^2\big(u^1\big)\big)
\\
\phantom{|d\mathcal{X}|^2=}{}
+\big(a_2-\big(a_2-a_0-(a_1-a_0)\sin^2\big(u^1\big)\big)\sin^2\big(u^2\big)\big)\big(du^2\big)^2.
\end{gather*}

Thus the condition $|d\mathcal{X}_z|^2=|d\mathcal{X}|^2$ becomes
\begin{gather*}
f^2\big(z,u^1\big)+\big(a_2-a_0-(a_1-a_0)\sin^2\big(u^1\big)\big)=\mathrm{const}=
\frac{a_2-h'^2(z,u^2)}{\sin^2(u^2)},
\\
f'^2\big(z,u^1\big)+f^2\big(z,u^1\big)g'^2\big(z,u^1\big)=a_1-(a_1-a_0)\sin^2\big(u^1\big),
\end{gather*}
from where we get
\begin{gather}
h\big(z_1,u^2\big):=\int_0^{u^2}\sqrt{a_2-(a_2-z_1a_0)\sin^2(t)}dt,\nonumber\\
f\big(z_1,u^1\big):=\sqrt{(1-z_1)a_0+(a_1-a_0)\sin^2(u^1)},\nonumber\\
g\big(z_1,u^1\big):=\int_0^{u^1}\frac{\sqrt{(1-z_1)a_0a_1+(a_1-a_0)z_1a_0\sin^2(t)}}
{(1-z_1)a_0+(a_1-a_0)\sin^2(t)}dt.\label{eq:fg}
\end{gather}
Note that
\begin{gather}
f\big(0,u^1\big)\cos\big(g\big(0,u^1\big)\big)=\sqrt{a_0}\cos\big(u^1\big),\qquad
f\big(0,u^1\big)\sin\big(g\big(0,u^1\big)\big)=\sqrt{a_1}\sin\big(u^1\big),\label{eq:fg0}
\end{gather}
(we assume simplif\/ications of the form
$\sqrt{a}\sqrt{b}\simeq\sqrt{ab}$ with $\sqrt{\cdot}$ having the
usual def\/inition
$\sqrt{re^{i\theta}}:=\sqrt{r}e^{\frac{i\theta}{2}}$, $r>0$,
$-\pi<\theta\le\pi$, since the possible signs are accounted by
symmetries in the principal planes for quadrics and disappear at
the level of the linear element for their deformations), so~$\mathcal{X}=\mathcal{X}_0$.

The coordinates $\mathbf{x}_0$, $\mathbf{x}_1$, $\mathbf{x}_2$ of
$\mathcal{X}_1$ satisfy (modulo a sign at the second formula)
Peterson's formulae:
\begin{gather}
\sqrt{(\mathbf{x}_0)^2+(\mathbf{x}_1)^2}=\sqrt{a_1-a_0}\cos\big (u^2\big) \sin\big(u^1\big),\nonumber\\
\tan^{-1}\left(\frac{\mathbf{x}_1}{\mathbf{x}_0}\right)=
\frac{\sqrt{a_0}}{\sqrt{a_1-a_0}}\tanh^{-1}\big(\cos\big(u^1\big)\big),\nonumber\\
\mathbf{x}_2=\int_0^{u^2}\sqrt{a_2-(a_2-a_0)\sin^2(t)}dt.\label{eq:pf2}
\end{gather}

More generally
\begin{gather}
h\big(z_1,u^2\big):=\int_0^{u^2}\sqrt{h'^2(t)-z_1\sin^2(t)}dt,\nonumber\\
f\big(z_1,u^1\big):=\sqrt{z_1+f^2(u^1)},\nonumber\\
g\big(z_1,u^1\big):=\int_0^{u^1}\frac{\sqrt{z_1(f'^2(t)+f^2(t)g'^2(t))+f^4(t)g'^2(t)}}
{z_1+f^2(t)}dt\label{eq:bla}
\end{gather}
give Peterson's $1$-dimensional family of deformations~(\ref{eq:surf}) with common conjugate system $(u^1,u^2)$.


\section{Peterson's deformations of higher dimensional quadrics}\label{sec:pet1}

Again we shall discuss only the case of quadrics with center and
having distinct eigenvalues of the quadratic part def\/ining the
quadric, without insisting on totally real cases and deformations
(when the linear elements are real valued).

\begin{remark}
A metric classif\/ication of all (totally real) quadrics in
$\mathbb{C}^{n+1}$ requires the notion of {\it symmetric Jordan}
canonical form of a symmetric complex matrix (see, e.g.~\cite{H}). The symmetric Jordan blocks are:
\begin{gather*}
J_1:=0=0_{1,1}\in\mathbf{M}_1(\mathbb{C}),\qquad
J_2:=f_1f_1^T\in\mathbf{M}_2(\mathbb{C}),\qquad
J_3:=f_1e_3^T+e_3f_1^T\in\mathbf{M}_3(\mathbb{C}),
\\
J_4:=f_1\bar f_2^T+f_2f_2^T+\bar
f_2f_1^T\in\mathbf{M}_4(\mathbb{C}),\qquad J_5:= f_1\bar
f_2^T+f_2e_5^T+e_5f_2^T+\bar
f_2f_1^T\in\mathbf{M}_5(\mathbb{C}),
\\
J_6:= f_1\bar f_2^T+f_2\bar f_3^T+f_3f_3^T+\bar f_3f_2^T+\bar
f_2f_1^T\in\mathbf{M}_6(\mathbb{C}),
\end{gather*}
etc., where $f_j:=\frac{e_{2j-1}-ie_{2j}}{\sqrt{2}}$ are the
standard isotropic vectors (at least the blocks $J_2$, $J_3$ were
known to the classical geometers). Any symmetric complex matrix
can be brought via conjugation with a complex rotation to the
symmetric Jordan canonical form, that is a matrix block
decomposition with blocks of the form $a_jI_p+J_p$; totally real
quadrics are obtained for eigenvalues $a_j$ of the quadratic part
def\/ining the quadric being real or coming in complex conjugate
pairs $a_j$, $\bar a_j$ with subjacent symmetric Jordan blocks of
same dimension $p$.
\end{remark}

Consider the quadric $\sum\limits_{j=0}^n\frac{(x_j)^2}{a_j}=1$, $
a_j\in\mathbb{C}^*$ distinct with parametrization given by the
spherical coordinates on the unit sphere
$\mathbb{S}^n\subset\mathbb{C}^{n+1}$
\[
\mathcal{X}=\sqrt{a_0}\mathbf{C}_0e_0+
\sum_{k=1}^n\sqrt{a_k}\mathbf{C}_k\sin\big(u^k\big)e_k,\qquad
\mathbf{C}_k:=\prod_{j=k+1}^n\cos\big(u^j\big).
\]
The correct generalization of (\ref{eq:surf}) allows us to build
Peterson's deformations of higher dimensional quadrics. With an
eye to the case $n=2$ we make the natural ansatz
\begin{gather}
\mathcal{X}_{\mathbf{z}}=\sum_{k=1}^{n-1}\mathbf{C}_kf_k\big(\mathbf{z},u^k\big)
\big(\cos\big(g_k\big(\mathbf{z},u^k\big)\big)e_{2k-2}
+\sin\big(g_k\big(\mathbf{z},u^k\big)\big)e_{2k-1}\big)+h\big(\mathbf{z},u^n\big)e_{2n-2}\label{eq:surf3}
\end{gather}
with the parameter(s) $\mathbf{z}=(z_1,z_2,\dots )$ to be determined
later.

We have
\begin{gather*}
|d\mathcal{X}_{\mathbf{z}}|^2=\sum_{k=1}^{n-1}\big[\mathbf{C}_k^2\big(f_k'^2\big(\mathbf{z},u^k\big)
+f_k^2\big(\mathbf{z},u^k\big)g_k'^2\big(\mathbf{z},u^k\big)\big)\big(du^k\big)^2
+\tfrac{1}{2}d\big(\mathbf{C}_k^2\big)d\big(f_k^2\big(\mathbf{z},u^k\big)\big)
\\
\phantom{|d\mathcal{X}_{\mathbf{z}}|^2=}{}
+f_k^2(\mathbf{z},u^k)(d\mathbf{C}_k)^2\big]
+h'^2\big(\mathbf{z},u^n\big)\big(du^n\big)^2,\\
|d\mathcal{X}|^2=a_0(d\mathbf{C}_0)^2
+\sum_{k=1}^na_k\big(d\big(\mathbf{C}_k\sin\big(u^k\big)\big)\big)^2.
\end{gather*}
Comparing the coef\/f\/icients of $(du^n)^2$ from
$|d\mathcal{X}_{\mathbf{z}}|^2=|d\mathcal{X}|^2$ we get
\begin{gather*}
\frac{1}{\cos^2(u^n)}\left[\mathbf{C}_1^2\big(f_1^2\big(\mathbf{z},u^1\big)-a_0-(a_1-a_0)\sin^2\big(u^1\big)\big)+
\sum_{k=2}^{n-1}\mathbf{C}_k^2\big(f_k^2\big(\mathbf{z},u^k\big)-a_k\sin^2\big(u^k\big)\big)\right]
\\
\qquad {}=\mathrm{const}=\frac{a_n\cos^2(u^n)-h'^2(\mathbf{z},u^n)}{\sin^2(u^n)}
\end{gather*}
from where we get with $z_0:=1$:
\begin{gather*}
f_k^2\big(z_{k-1},z_k,u^k\big):=(z_{k-1}-z_k)a_0+(a_k-z_{k-1}a_0)\sin^2\big(u^k\big),\qquad
k=1,\dots ,n-1,
\\
h'^2\big(z_{n-1},u^n\big):=a_n-(a_n-z_{n-1}a_0)\sin^2\big(u^n\big).
\end{gather*}
Now we have
\begin{gather*}
(d\mathbf{C}_0)^2=\sum_{k=1}^{n-1}\big[z_{k-1}(d\mathbf{C}_{k-1})^2-z_k(d\mathbf{C}_k)^2\big]
+z_{n-1}(d\mathbf{C}_{n-1})^2=\sum_{k=1}^{n-1}\!\big[z_{k-1} \big(\mathbf{C}_k^2\sin^2\big(u^k\big)\big(du^k\big)^2
\\
\phantom{(d\mathbf{C}_0)^2=}{}
-\tfrac{1}{2}d\big(\mathbf{C}_k^2\big)d\big(\sin^2(u^k)\big)+\cos^2\big(u^k\big)\big(d\mathbf{C}_k\big)^2\big)
-z_k\big(d\mathbf{C}_k\big)^2\big] +z_{n-1}\big(d\mathbf{C}_{n-1}\big)^2,
\\
\big(d\big(\mathbf{C}_k\sin\big(u^k\big)\big)\big)^2=\mathbf{C}_k^2\cos^2\big(u^k\big)\big(du^k\big)^2
+\tfrac{1}{2}d\big(\mathbf{C}_k^2\big)d\big(\sin^2\big(u^k\big)\big)+\sin^2\big(u^k\big)\big(d\mathbf{C}_k\big)^2,
\end{gather*}
so
\begin{gather*}
|d\mathcal{X}|^2=\sum_{k=1}^{n-1}\big[\mathbf{C}_k^2\big(a_k-
(a_k-z_{k-1}a_0)\sin^2\big(u^k\big)\big)\big(du^k\big)^2
+\tfrac{1}{2}(a_k-z_{k-1}a_0)d\big(\mathbf{C}_k^2\big)d\big(\sin^2\big(u^k\big)\big)
\\
\phantom{|d\mathcal{X}|^2=}{}
+\big((z_{k-1}-z_k)a_0+
(a_k-z_{k-1}a_0)\sin^2\big(u^k\big)\big)\big(d\mathbf{C}_k\big)^2\big]\\
\phantom{|d\mathcal{X}|^2=}{}
+\big(a_n-(a_n-z_{n-1}a_0)\sin^2\big(u^n\big)\big)\big(du^n\big)^2,
\\
0=|d\mathcal{X}_{\mathbf{z}}|^2-|d\mathcal{X}|^2=
\sum_{k=1}^{n-1}\mathbf{C}_k^2\big(f_k'^2\big(\mathbf{z},u^k\big)
+f_k^2\big(\mathbf{z},u^k\big)g_k'^2\big(\mathbf{z},u^k\big)\\
\phantom{0=}{}-a_k+
(a_k-z_{k-1}a_0)\sin^2\big(u^k\big)\big)\big(du^k\big)^2,
\end{gather*}
so we f\/inally get (\ref{eq:fgn}).

For $z_1=z_2=\dots =z_{n-1}=0$ we get $g_2=\dots =g_{n-1}=0$ and using
(\ref{eq:fg0}) we get $\mathcal{X}=\mathcal{X}_{\mathbf{0}}$ with
$\mathbb{C}^{n+1}\hookrightarrow\mathbb{C}^{2n-1}$ as
$(x_0,x_1,\dots ,x_n)\mapsto(x_0,x_1,x_2,0,x_3,0,\dots ,x_{n-1},0,x_n)$.

For $z_1=z_2=\dots =z_{n-1}=1$ we get
$\mathcal{X}_{\mathbf{1}}=(\mathbf{x}_0,\dots ,\mathbf{x}_{2n-2})$
given by Peterson's formulae (\ref{eq:pfn}).

More generally and with $z_0:=0$
\begin{gather}
f_k\big(z_{k-1},z_k,u^k\big):=\sqrt{z_k+f_k^2(u^k)-z_{k-1}\cos^2(u^k)},\qquad k=1,\dots ,n-1,\nonumber\\
g_k\big(z_{k-1},z_k,u^k\big):=\int_0^{u^k}\frac{\sqrt{f_k'^2(t)+f_k^2(t)g_k'^2(t)-
(f_k'^2(z_{k-1},z_k,t)+z_{k-1}\sin^2(t))}}
{f_k(z_{k-1},z_k,t)}dt,\nonumber\\
h\big(z_{n-1},u^n\big):=\int_0^{u^n}\sqrt{h'^2(t)-z_{n-1}\sin^2(t)}dt \label{eq:fgr}
\end{gather}
give an $(n-1)$-dimensional family of deformations
(\ref{eq:surf3}); for $g_k(u^k)=0$, $k=2,\dots ,n-1$ we have
$\mathcal{X}_{\mathbf{0}}\subset\mathbb{C}^{n+1}$.

\section{The common conjugate system\\ and non-degenerate joined second fundamental forms}\label{sec:pet2}

The fact that $(u^1,\dots ,u^n)$ is a conjugate system on
$\mathcal{X}_{\mathbf{0}}$ is clear since we have
\[
\partial_{u^k}\partial_{u^j}\mathcal{X}_\mathbf{0}=
-\tan\big(u^j\big)\partial_{u^k}\mathcal{X}_\mathbf{0},\qquad 1\le k<j\le n.
\]
With the normal f\/ield
\[
\hat
N_{\mathbf{0}}:=(\sqrt{a_0})^{-1}\mathbf{C}_0e_0+
\sum_{k=1}^n(\sqrt{a_k})^{-1}\mathbf{C}_k\sin\big(u^k\big)e_k
\]
we have $\hat
N_{\mathbf{0}}^Td^2\mathcal{X}_{\mathbf{0}}=-\sum\limits_{k=1}^n\mathbf{C}_k^2(du^k)^2$.
To see that $(u^1,\dots ,u^n)$ is a conjugate system on
\begin{gather*}
\mathcal{X}=(x_0,\dots ,x_{2n-2}):=\sum_{k=1}^{n-1}\mathbf{C}_kf_k\big(u^k\big)\big(\cos\big(g_k\big(u^k\big)\big)e_{2k-2}
+\sin\big(g_k\big(u^k\big)\big)e_{2k-1}\big)\\
\phantom{\mathcal{X}=(x_0,\dots ,x_{2n-2}):=}{} +h\big(u^n\big)e_{2n-2}
\end{gather*}
we have again
$\partial_{u^k}\partial_{u^j}\mathcal{X}=-\tan(u^j)\partial_{u^k}\mathcal{X}$, $1\le k<j\le n$; again the $n-1$ f\/ields
\begin{gather*}
\partial_{u^1}\mathcal{X}|_{u^1,u^2,\dots ,\widehat{u^k},\dots ,u^n=\mathrm{const}},\qquad
\partial_{u^2}\mathcal{X}|_{u^1,u^2,\dots ,\widehat{u^k},\dots ,u^n=\mathrm{const}},\qquad
\dots ,\\
\widehat{\partial_{u^k}\mathcal{X}}|_{u^1,u^2,\dots ,\widehat{u^k},\dots ,u^n=\mathrm{const}},\qquad
\dots ,
\qquad
\partial_{u^n}\mathcal{X}|_{u^1,u^2,\dots ,\widehat{u^k},\dots ,u^n=\mathrm{const}},\qquad
k=1,\dots ,n
\end{gather*}
generate ruled $n$-dimensional developables in $\mathbb{C}^{2n-1}$
because the only term producing shape is
$\partial_{u^k}\partial_{u^k}\mathcal{X}$.

For the non-degenerate joined second fundamental forms property we
have
\begin{gather*}
u^n=h^{-1}(x_{2n-2}),\qquad h'\big(u^n\big)du^n=dx_{2n-2},\qquad
u^k=g_k^{-1}\left(\tan^{-1}\left(\frac{x_{2k-1}}{x_{2k-2}}\right)\right),
\\
\mathbf{C}_k^2f_k^2(u^k)g'_k(u^k)du^k=x_{2k-2}dx_{2k-1}-x_{2k-1}dx_{2k-2},
\qquad k=1,\dots ,n-1
\end{gather*}
and $\mathcal{X}$ is given implicitly by the zeroes of the
functionally independent
\[
F_k:=(x_{2k-2})^2+(x_{2k-1})^2-\mathbf{C}_k^2f_k^2\big(u^k\big),\qquad
k=1,\dots ,n-1.
\]
We have the natural linearly independent normal f\/ields
\begin{gather*}
N_k:=\nabla F_k=2(x_{2k-2}e_{2k-2}+x_{2k-1}e_{2k-1})
-\frac{2f_k'(u^k)(-x_{2k-1}e_{2k-2}+x_{2k-2}e_{2k-1})}{f_k(u^k)g'_k(u^k)}
\\
\phantom{N_k:=}{} +2\mathbf{C}_k^2f_k^2(u^k)\left[\sum_{j=k+1}^{n-1}\frac{\tan(u^j)
(-x_{2j-1}e_{2j-2}+x_{2j-2}e_{2j-1})}
{\mathbf{C}_j^2f_j^2(u^j)g'_j(u^j)}+\frac{\tan(u^n)e_{2n-2}}{h'(u^n)}\right],\\
\phantom{N_k:=}{} k=1,\dots ,n-1,
\end{gather*} and
\begin{gather*}
\partial_{u^l}\partial_{u^l}\mathcal{X}=-\sum_{j=1}^{l-1}(x_{2j-2}e_{2j-2}+x_{2j-1}e_{2j-1})
+\left(\frac{f_l''(u^l)}{f_l(u^l)}-g_l'^2\big(u^l\big)\right)(x_{2l-2}e_{2l-2}+x_{2l-1}e_{2l-1})
\\
\phantom{\partial_{u^l}\partial_{u^l}\mathcal{X}=}{} +g'_l\big(u^l\big)\left(\frac{2f'_l(u^l)}{f_l(u^l)}+\frac{g_l''(u^l)}{g_l'(u^l)}\right)
(-x_{2l-1}e_{2l-2}+x_{2l-2}e_{2l-1}),\qquad l=1,\dots ,n-1,
\\
\partial_{u^n}\partial_{u^n}\mathcal{X}=-\sum_{l=1}^{n-1}(x_{2l-2}e_{2l-2}+x_{2l-1}e_{2l-1})
+h''\big(u^n\big)e_{2n-2},
\end{gather*}
and the second fundamental form
\begin{gather*}
N_k^Td^2\mathcal{X}=2\mathbf{C}_k^2f_k^2\Bigg[\left(\frac{f_k''(u^k)}{f_k(u^k)}-g_k'^2\big(u^k\big)
-\frac{f_k'(u^k)}{f_k(u^k)}\left(\frac{2f'_k(u^k)}{f_k(u^k)}+
\frac{g_k''(u^k)}{g_k'(u^k)}\right)\right)\big(du^k\big)^2
\\
\phantom{N_k^Td^2\mathcal{X}=}{}
+\sum_{l=k+1}^{n-1}\left(\tan\big(u^l\big)\left(\frac{2f'_l(u^l)}{f_l(u^l)}
+\frac{g_l''(u^l)}{g_l'(u^l)}\right)-1\right)\big(du^l\big)^2\\
\phantom{N_k^Td^2\mathcal{X}=}{}
+\left(\frac{\tan(u^n)h''(u^n)}{h'(u^n)}-1\right)\big(du^n\big)^2\Bigg],\qquad k=1,\dots ,n-1.
\end{gather*}

For Peterson's deformations of higher dimensional quadrics we have
\begin{gather*}
N_k^Td^2\mathcal{X}=
-2a_0\mathbf{C}_k^2f_k^2\left(\frac{a_kz_{k-1}(du^k)^2}{g_k'^2(u^k)f_k^4(u^k)}
+\sum_{l=k+1}^{n-1}\frac{a_l(z_{l-1}-z_l)(du^l)^2}{g_l'^2(u^l)f_l^4(u^l)}
+\frac{a_n(du^n)^2}{a_0h'^2(u^n)}\right).
\end{gather*}
It is now enough to check the open non-degenerate joined second
fundamental forms property only for $\mathbf{z}=(1,1,\dots ,1)$. Thus
with $\delta:=\frac{a_n}{a_0\sin^2(u^n)+a_n\cos^2(u^n)}$ we need
\begin{gather*}
0\neq\begin{vmatrix}
\mathbf{C}_1&\mathbf{C}_2&\mathbf{C}_3&\dots &\mathbf{C}_{n-1}&\delta^{-1}\mathbf{C}_n\\
\frac{a_1}{a_1-a_0}&0&0&\dots &0&\sin^2\big(u^1\big)\\
0&\frac{a_2}{a_2-a_0}&0&\dots &0&\sin^2\big(u^2\big)\\
0&0&\frac{a_3}{a_3-a_0}&\dots &0&\sin^2\big(u^3\big)\\
\vdots&\vdots&\vdots&\cdots&\vdots&\vdots\\
0&0&0&\dots &\frac{a_{n-1}}{a_{n-1}-a_0}&\sin^2\big(u^{n-1}\big)\end{vmatrix}
\end{gather*}
almost everywhere, which is straightforward.

\begin{remark}
Note that a-priori $\mathcal{X}_{\mathbf{1}}$ comes close to lie
in a degenerate deformation of $\mathbb{C}^{n+1}$ in
$\mathbb{C}^{2n-1}$:  $\hat
N_{\mathbf{0}}^Td^2\mathcal{X}_{\mathbf{0}}-
\big(\sum\limits_{k=1}^{n-1}\frac{1}{2a_k}N_k\big)^Td^2\mathcal{X}_{\mathbf{1}}$
depends only on $(du^n)^2$ and this is as closest to $0$ as we can
get.
\end{remark}
